\section{Introdução}
\label{section:intro}

Tecnologias inovadoras baseadas em DNA, que buscam a caracterização da diversidade biológica em larga escala, têm sido amplamente difundidas. Uma das técnicas mais utilizadas é a análise de amostras de DNA ambiental (eDNA) \cite{Ficetola.etal2008}. As amostras de eDNA são amplificadas utilizando iniciadores (primers) para regiões de interesse previamente definidas, e posteriormente submetida ao sequenciamento de alta geração de dados (High Throughput Sequencing – HTS). Esse processo garante a identificação taxonômica de um pool de espécies a partir de uma pequena amostra, sem a necessidade de manuseio e captura dos organismos de interesse, metodologia conhecida como eDNA metabarcoding \cite{Cristescu2014,Deiner.etal2017,Fahner.etal2016}.

Os primeiros estudos para detecção de vertebrados por eDNA foram realizados por Ficetola et al., 2008, que observaram a presença de \textit{Lithobates catesbeianus} (rã-touro-americana) em ambientes naturais e controlados. Para peixes, vários estudos mostraram eficiência de detecção nos diversos ecossistemas, como lagoas, rios, riachos e em regiões marítimas \cite{Takahara.etal2013,Sigsgaard.etal2015,Minamoto.etal2012,Thomsen.etal2012,Kelly.etal2014}. Nos ambientes aquáticos, o material genético detectado é originário de diversas fontes como tecidos deteriorados ou células da pele e estão depositados na camada superficial da água. 

Estudos com eDNA metabarcoding são ferramentas poderosas que minimizam o contato com o objeto de estudo e contribuem significativamente para o monitoramento da biodiversidade, gerando conhecimento no intuito de diminuir os impactos sobre o meio ambiente. A metodologia de eDNA metabarcoding vem sendo aplicada em projetos de monitoramentos ambiental, no intuito de avaliar a eficiência de estratégias de manejo, bem como em estudos sobre conservação da biodiversidade, atribuindo resultados rápidos e confiáveis para um conjunto de decisões regulatórias e de conservação \cite{Oliveira.etal2023}.

Porém, para análise de composição de comunidades e avaliação de biodiversidade com a metodologia de DNA metabarcoding são necessários protocolos eficientes em cada uma de suas etapas. Uma das etapas mais críticas  desses estudos é a seleção do marcador molecular e dos iniciadores, que determina a capacidade de detecção da comunidade alvo. Para o monitoramento de comunidades ictiológicas, o gene mitocondrial 12S rRNA é amplamente consolidado como marcador padrão para DNA ambiental.

Os primers MiFish \cite{Miya.etal2015a} são um conjunto de iniciadores universais projetados especificamente para amplificar um fragmento curto de aproximadamente 170 pares de base do gene 12S rRNA. Esses primers demonstraram uma alta capacidade de discriminar táxons de peixes, tanto teleósteos quanto elasmobrânquios. Para a validação do MiFish, os autores originais realizaram análises que contrastaram ambientes naturais com ambientes controlados cuja composição da biodiversidade era previamente conhecida, a fim de testar e aferir a sensibilidade dos iniciadores.

Diante disso, este projeto utilizou uma fração dos dados de sequenciamento depositados no DNA Databank of Japan (DDBJ) pelo bioprojeto original do MiFish. O objetivo geral de nosso projeto foi avaliar e comparar a estrutura, riqueza e composição taxonômica da comunidade de peixes entre o ecossistema natural do Recife de Bise e o ambiente controlado do Aquário Churaumi utilizando metabarcoding de eDNA. Para alcançar este propósito, os objetivos específicos delineados foram a implementação de um pipeline reprodutível no ambiente QIIME 2 para a resolução de Variantes de Sequência de Amplicons (ASVs), a comparação dos índices de diversidade alfa (riqueza observada) e beta (similaridade por distância de Jaccard) entre os dois ambientes, e, por fim, a avaliação comparativa do perfil e composição das famílias de peixes detectadas em cada ambiente.